\subsection{Sélection physique comme relation naturelle}
En dehors du centre électronique, une espèce ne doit pas choisir par décret une cellule unique. Les douze multisections non centrales sont des fibres stables sous l’inversion cellulaire et n’admettent généralement pas de section ponctuelle équivariante. La réalisation correcte commence par une relation

\[
\mathscr S\subset \mathcal P\times\mathcal R_{324},
\qquad
\mathscr S(X)=\{\gamma:(X,\gamma)\in\mathscr S\},
\]
dont la valeur peut être un point, une orbite ou une fibre. Pour l’électron, \(\mathscr S(e)=\{\gamma_e\}\) est de rang un parce que \((5,5)\) est l’unique point fixe. Pour la couleur, la saveur et la génération, \(\mathscr S(X)\) conserve l’orbite correspondante jusqu’à ce qu’un couplage physique la comprime.

La relation est admissible si elle satisfait simultanément les conditions suivantes :

\begin{enumerate}
\item elle dépend uniquement de la représentation, de l’orientation, de la mémoire, de la frontière, de la charge, de la saveur, de la couleur et de la génération construites avant la comparaison ;

\item elle commute avec la conjugaison de particule et d’antiparticule ;
\item elle est équivariante relativement à l’action chromatique et à l’inversion cellulaire ;
\item elle est compatible avec le raffinement, le retour nonadique et la lecture dodécaphasique ;
\item elle conserve toute multiplicité qui n’a pas été éliminée par un couplage ;
\item elle reste invariante lorsqu’on permute ou remplace les valeurs externes de masse.
\end{enumerate}
Le sixième point est le contrôle décisif. Si la modification de la colonne expérimentale change \(\mathscr S(X)\), la construction est rétrospective. Si la relation demeure et que seul le résidu de comparaison change, la direction causale est préservée.

\Needspace{7\baselineskip}
